% ----------------------------------------------------------
% Desenvolvimento
% ----------------------------------------------------------
\chapter{Desenvolvimento}
\label{cap:desenvolvimento}
% ----------------------------------------------------------

A implementação da metodologia de análise de redes complexas aplicada a competições olímpicas demandou pipeline técnico abrangente, desde curadoria de dados históricos até desenvolvimento de ferramentas interativas de visualização. O processo envolveu decisões arquiteturais sobre estruturas de dados, seleção de bibliotecas especializadas, e parametrização de algoritmos de grafos direcionados ponderados. A escolha de linguagem Python e seu ecossistema de ciência de dados viabilizou integração eficiente entre manipulação de dados tabulares, análise de redes e construção de interfaces web interativas.

\section{Seleção e Preparação dos Dados}

Com base na fundamentação teórica desenvolvida no Capítulo~\ref{cap:revisao}, foram selecionadas três modalidades que representam os tipos fundamentais de interdependência competitiva: Swimming (esporte misto com provas individuais e revezamentos, permitindo análise comparativa dentro do mesmo contexto \citeonline{Silva2019}), Basketball e Football (esportes coletivos genuínos oferecendo contraste em dinâmica de jogo \citeonline{TeamSportsNetwork2017}). Esta seleção garante cobertura dos tipos fundamentais de dinâmicas competitivas identificados na taxonomia de interdependência competitiva (Seção~\ref{sec:taxonomia}).

A construção das redes competitivas inicia-se com curadoria de dados de dois conjuntos históricos olímpicos complementares. O dataset base \citeonline{Griffin2018OlympicDataset} abrange o período de Atenas 1896 até Rio de Janeiro 2016, totalizando 271.116 registros disponibilizados na plataforma Kaggle. Para estender a cobertura temporal, este conjunto foi integrado com dados de Tokyo 2020 \citeonline{Piterfm2021TokyoOlympics}, adicionando 2.411 registros da olimpíada realizada em 2021. O processo de integração envolveu normalização de formatos de nomes de eventos (adequando nomenclatura de Basketball e Football ao padrão histórico), mapeamento correto de campos de gênero (convertendo codificação M/W/X para formato M/F padrão), e remoção de registros duplicados por ID de atleta para garantir consistência. O dataset consolidado final abrange 125 anos de história olímpica (1896-2021), totalizando 273.527 registros que correspondem a 135.571 atletas únicos competindo em 66 modalidades esportivas distribuídas em 765 eventos distintos, representando 230 comitês olímpicos nacionais ao longo de 53 edições dos jogos.

O dataset estrutura-se em 15 colunas que capturam informações demográficas, bioantropométricas e competitivas dos atletas (Tabela~\ref{tab:dataset_metadata}): ID (identificador único do atleta), Name (nome completo), Sex (gênero: M/F), Age (idade no momento da competição), Height (altura em centímetros), Weight (peso em quilogramas), Team (país representado), NOC (código do comitê olímpico nacional em formato ISO de 3 caracteres), Games (identificador da edição olímpica), Year (ano dos jogos), Season (temporada: Summer/Winter), City (cidade sede), Sport (modalidade esportiva), Event (evento específico dentro da modalidade), e Medal (tipo de medalha conquistada: Gold/Silver/Bronze ou NA para não medalhistas). A presença do campo ID permite rastreamento longitudinal de atletas que participaram de múltiplas edições olímpicas, viabilizando análises de carreiras competitivas completas.

% Tabela de Metadados do Dataset Olímpico
\begin{table}[htbp]
\centering
\caption{Estrutura de metadados do dataset olímpico histórico}
\label{tab:dataset_metadata}
\footnotesize
\begin{tabular}{lp{2.5cm}p{8cm}}
\toprule
\textbf{Coluna} & \textbf{Tipo} & \textbf{Descrição} \\
\midrule
ID & Inteiro & Identificador único do atleta. Permite rastreamento longitudinal de participações em múltiplas edições olímpicas. \\
\midrule
Name & String & Nome completo do atleta conforme registrado nos arquivos olímpicos oficiais. \\
\midrule
Sex & Categórico (M/F) & Gênero do atleta: Masculino (M) ou Feminino (F). \\
\midrule
Age & Inteiro & Idade do atleta no momento da competição. Valores faltantes (NA) presentes para competidores históricos. \\
\midrule
Height & Numérico & Altura do atleta em centímetros. Valores faltantes (NA) comuns em registros anteriores a 1960. \\
\midrule
Weight & Numérico & Peso do atleta em quilogramas. Valores faltantes (NA) comuns em registros anteriores a 1960. \\
\midrule
Team & String & Nome completo da equipe/país representado pelo atleta (ex: "United States", "Brazil"). \\
\midrule
NOC & String (3 letras) & Código do Comitê Olímpico Nacional em formato padrão ISO (ex: "USA", "BRA", "CHN"). Campo mais confiável que Team para análises por país. \\
\midrule
Games & String & Identificador da edição olímpica no formato "Ano Temporada" (ex: "2016 Summer", "2014 Winter"). \\
\midrule
Year & Inteiro & Ano de realização dos jogos olímpicos (1896 a 2021). \\
\midrule
Season & Categórico & Temporada da competição: Verão (Summer) ou Inverno (Winter). \\
\midrule
City & String & Cidade sede dos jogos olímpicos (ex: "Athens", "Rio de Janeiro", "Beijing"). \\
\midrule
Sport & String & Modalidade esportiva (ex: "Swimming", "Basketball", "Football", "Athletics"). \\
\midrule
Event & String & Evento específico dentro da modalidade (ex: "Swimming Men's 100 metres Freestyle"). \\
\midrule
Medal & Categórico & Tipo de medalha conquistada: Gold, Silver, Bronze, ou NA para participações sem pódio. \\
\bottomrule
\end{tabular}
\fonte{\citeonline{Griffin2018OlympicDataset}}
\end{table}


Do total de 271.116 registros, apenas 39.783 (14.7\%) correspondem a conquistas de medalhas, distribuídos em 13.372 medalhas de ouro, 13.116 de prata e 13.295 de bronze. Para o presente trabalho, foram aplicados filtros sistemáticos que resultaram em conjunto reduzido mas representativo: seleção exclusiva de medalhistas (39.783 registros), restrição a jogos olímpicos de verão (excluindo jogos de inverno), e foco em três modalidades esportivas estrategicamente escolhidas (Swimming, Basketball, Football). Adicionalmente, os dados foram segregados por gênero, gerando redes independentes para competições masculinas e femininas. Este processo de filtragem resultou em conjunto final de 4.659 atletas medalhistas distribuídos nas três modalidades selecionadas, constituindo os vértices das redes construídas para análise.

\subsection{Limpeza e Integridade dos Dados}

Para garantir integridade estrutural das redes competitivas, foram aplicados filtros sistemáticos de limpeza de dados durante o pipeline de processamento. O processo de limpeza envolveu duas etapas críticas que asseguram a premissa fundamental da modelagem: redes competitivas baseiam-se em compartilhamento de pódio, exigindo no mínimo dois atletas diferentes em cada evento para criação de arestas.

Primeiro, foram identificados e removidos \textbf{2 registros duplicados} correspondentes a atletas com múltiplas entradas idênticas no mesmo evento (mesma medalha, mesmo ano, mesmo evento, mesmo sexo). Estes duplicados, principalmente em eventos de vela (Sailing) de edições antigas (1900), representavam inconsistências de registro histórico que causariam sobreponderação artificial de conexões competitivas.

Segundo, foram analisados os \textbf{662 pódios válidos} dos três esportes selecionados. O processo de filtragem opera em dois pontos do pipeline: durante carregamento inicial dos dados (remoção preventiva de duplicatas e pódios com menos de 2 atletas únicos por evento) e durante construção das redes (validação adicional assegurando que cada evento processado contém pelo menos 2 atletas diferentes antes de criar arestas). Esta abordagem em camadas garante robustez estrutural, eliminando possibilidade de nós isolados decorrentes de inconsistências de registro histórico.

O conjunto final de \textbf{4.659 atletas medalhistas} distribuídos em \textbf{662 pódios compartilhados} nas três modalidades selecionadas constitui base limpa e estruturalmente íntegra para modelagem de redes competitivas, assegurando que todo atleta incluído participou de pelo menos uma interação competitiva direta (compartilhamento de pódio) modelável como aresta direcionada ponderada. Esta limpeza resultou em redes sem nós isolados, característica fundamental para cálculo correto de métricas de centralidade e detecção confiável de comunidades.

Cada vértice nas redes competitivas representa atleta único identificado pelo ID, com atributos extraídos do dataset original preservados: sexo, idade, altura, peso, equipe, código nacional olímpico, jogos, cidade e medalha. Este enriquecimento atributivo permite análises multifacetadas que correlacionam estrutura de rede com características demográficas, geográficas e temporais \citeonline{Sousa2018Computational}. O processo de construção segrega os dados por esporte e gênero, gerando grafos direcionados independentes para cada subconjunto \citeonline{Radicchi2016OlympicDominance}.

As arestas capturam a essência competitiva: em cada evento, medalhistas são ordenados hierarquicamente (Ouro > Prata > Bronze), estabelecendo arestas direcionadas do medalhista de posição inferior para o superior. O peso da aresta reflete disparidade de medalhas: Ouro sobre Bronze (peso 5), Ouro sobre Prata (peso 3), Prata sobre Bronze (peso 2). Pesos são acumulados em interações múltiplas, refletindo rivalidades persistentes. Esta ponderação hierárquica quantifica dominância competitiva \citeonline{Leicht2008CommunityStructure}. Os grafos resultantes exibem heterogeneidade, com número de vértices e arestas variando por esporte e gênero.

\section{Implementação Técnica e Pipeline de Análise}

A implementação do pipeline de análise foi desenvolvida em Python, linguagem amplamente utilizada em ciência de dados pela riqueza de bibliotecas especializadas e flexibilidade para manipulação de estruturas complexas. O ecossistema Python oferece ferramentas consolidadas para análise de redes, processamento de dados e visualização interativa, facilitando a reprodutibilidade e escalabilidade do projeto. O projeto integra bibliotecas especializadas em diferentes etapas do pipeline: NetworkX (criação, manipulação e análise de redes complexas, fornecendo estruturas eficientes para grafos direcionados e ponderados), Pandas e NumPy (processamento de dados tabulares e cálculos numéricos eficientes), Python-Louvain (implementação do algoritmo de Louvain para detecção de comunidades), e Plotly e Streamlit (visualizações interativas e construção de dashboards web).

O fluxo de processamento segue arquitetura modular dividida em três componentes principais. O módulo de mineração de dados realiza extração, transformação e carga dos dados olímpicos, filtrando medalhistas e segregando por esporte e gênero. O módulo de análise de redes implementa algoritmos de teoria de grafos para cálculo de métricas de centralidade, detecção de comunidades e caracterização estrutural, operando sobre os grafos construídos e gerando métricas numéricas para cada vértice. O módulo de visualização consiste em dashboard interativo que integra os resultados, permitindo exploração filtrada por esporte, gênero, período temporal e outras dimensões através de visualizações comparativas e tabelas detalhadas.

\section{Métricas de Análise de Redes}

PageRank, originalmente desenvolvido para ranking de páginas web, foi adaptado para contextos de redes competitivas esportivas \citeonline{Shiue2025ApplyingPageRank}. O algoritmo atribui importância a cada vértice baseado não apenas no número de conexões recebidas, mas também na importância dos vértices que o conectam. No contexto olímpico, atletas que derrotaram competidores já reconhecidos recebem pontuações mais elevadas, criando hierarquia de prestígio competitivo. A métrica calcula a probabilidade estacionária de um passeio aleatório na rede, ponderado pelos pesos das arestas. Atletas centrais na estrutura competitiva emergem naturalmente através deste processo iterativo, independentemente de métricas tradicionais como contagem absoluta de medalhas.

Além do PageRank, foram calculadas métricas complementares de centralidade \citeonline{Freeman1979Centrality}. Centralidade de grau quantifica conexões diretas de entrada (in-degree) e saída (out-degree), onde in-degree elevado indica atletas frequentemente superados por outros, enquanto out-degree alto representa aqueles que superaram muitos competidores. Centralidade de intermediação (betweenness) identifica atletas que atuam como pontes entre diferentes grupos competitivos, com valores elevados sugerindo posição estratégica na rede conectando comunidades ou eras distintas. Centralidade de proximidade (closeness) mede distância média de um vértice a todos os demais, identificando atletas centralmente posicionados na estrutura global da rede.

A identificação de comunidades nas redes competitivas foi realizada através do algoritmo de Louvain \citeonline{Blondel2008Louvain}, método hierárquico baseado em otimização de modularidade. A modularidade quantifica a qualidade de uma partição da rede, comparando a densidade de arestas dentro das comunidades com o esperado em uma rede aleatória \citeonline{Fortunato2010CommunityDetection}. O algoritmo opera em duas fases iterativas: na primeira, cada vértice é atribuído à comunidade que maximiza o ganho local de modularidade; na segunda, as comunidades identificadas são agregadas em super-nós, e o processo reinicia. Esta abordagem multi-escala revela estruturas hierárquicas presentes nas redes, com complexidade computacional eficiente mesmo para redes grandes. A detecção de comunidades permite identificar agrupamentos naturais de atletas baseados em padrões competitivos, revelando estruturas não evidentes em análises tradicionais.

Para cada comunidade detectada, foram calculadas métricas agregadas que permitem caracterização multidimensional: métricas temporais (período de atividade, entropia de Shannon da distribuição por décadas \citeonline{Rajaram2017ShannonEntropy}, era olímpica dominante), métricas geográficas (número de países representados, índice de concentração HHI \citeonline{Hirschman1964NationalPower}, entropia geográfica), e métricas de performance (distribuição de medalhas, estatísticas de PageRank, identificação de atletas de alta performance e atletas-ponte). Esta caracterização multi-facetada permite análise comparativa entre comunidades, revelando padrões estruturais distintos entre modalidades esportivas e contextos competitivos.

Além das métricas temporais, geográficas e de performance descritas, foram implementadas três análises complementares para caracterização multidimensional das comunidades detectadas, conforme fundamentação apresentada no Capítulo~\ref{cap:revisao}.

Para quantificar o nível de dominância competitiva de cada comunidade, calcula-se índice ponderado baseado na distribuição de medalhas utilizando sistema olímpico histórico (3 pontos para ouro, 2 para prata, 1 para bronze), fundamentado na sequência de Fibonacci estabelecida pela literatura em ranking olímpico \citeonline{deSousa2012NetworkBased}:

\begin{equation}
D = \frac{3 \times n_{\mathrm{ouro}} + 2 \times n_{\mathrm{prata}} + 1 \times n_{\mathrm{bronze}}}{n_{\mathrm{total}}}
\end{equation}

\noindent onde $n_{\mathrm{ouro}}$, $n_{\mathrm{prata}}$ e $n_{\mathrm{bronze}}$ representam contagens de medalhas de cada tipo, e $n_{\mathrm{total}}$ o número total de atletas medalhistas na comunidade. O índice varia entre 1,0 (exclusivamente bronze) e 3,0 (exclusivamente ouro). Este sistema de pontuação (3-2-1) difere da ponderação de arestas (5-3-2) empregada na construção das redes: enquanto pesos de arestas modelam distância competitiva entre atletas em confrontos diretos (maior peso = maior gap de performance), o índice de dominância quantifica valor hierárquico das medalhas para avaliação de perfil competitivo. Comunidades são classificadas empiricamente em três perfis: Elite ($D \geq 2,2$), Competitiva ($1,8 \leq D < 2,2$), e Participante ($D < 1,8$), baseado na distribuição observada do índice.

Para avaliar segregação estrutural, cada aresta da rede é classificada como intra-comunidade (conectando atletas da mesma comunidade) ou inter-comunidade (conectando atletas de comunidades distintas) \citeonline{Fortunato2010CommunityDetection}. A razão entre conexões internas e externas quantifica grau de segregação \citeonline{Leicht2008CommunityStructure}:

\begin{equation}
S = \frac{E_{\mathrm{intra}}}{E_{\mathrm{inter}} + \epsilon}
\end{equation}

\noindent onde $E_{\mathrm{intra}}$ e $E_{\mathrm{inter}}$ representam contagens de arestas intra e inter-comunidade, e $\epsilon = 10^{-6}$ evita divisão por zero. Valores elevados de $S$ indicam comunidades fortemente segregadas; valores próximos a zero caracterizam estruturas altamente interconectadas. Pares de comunidades com contagem excepcional de conexões mútuas são identificados como rivalidades estruturais.

Para identificar estrutura hierárquica no conjunto de comunidades, calcula-se PageRank médio de cada comunidade e estabelece-se classificação baseada em percentis \citeonline{Eriksson2022NetworkMetadata}. Comunidades são categorizadas em três níveis: Núcleo (percentil $\geq 80$, estruturalmente centrais), Intermediária ($40 \leq \mathrm{percentil} < 80$, conectoras), e Periférica (percentil $< 40$, marginalmente conectadas). Esta classificação revela que centralidade estrutural não necessariamente correlaciona com volume absoluto de medalhas, identificando comunidades em posições estratégicas independentemente de tamanho.

\section{Dashboard Interativo de Visualização}

Para facilitar exploração e interpretação dos resultados, foi desenvolvido dashboard interativo utilizando framework Streamlit, biblioteca Python para criação rápida de aplicações web analíticas. O dashboard integra todos os artefatos gerados pelo pipeline de análise: arquivos CSV consolidados contendo métricas de atletas e comunidades, arquivos JSON com resumos estatísticos por esporte e gênero, e dados de conectividade inter-comunidade para análise de rivalidades estruturais.

A interface organiza-se em abas temáticas que segmentam a análise em dimensões complementares. A aba de visão geral apresenta estatísticas descritivas básicas das redes construídas (número de atletas, arestas, densidade, modularidade), permitindo comparação visual entre modalidades e gêneros. As abas de centralidade fornecem rankings interativos baseados em PageRank, betweenness e degree centrality, com visualizações de distribuições e identificação de outliers estruturais. A aba de comunidades oferece caracterização detalhada dos agrupamentos detectados, incluindo perfis competitivos (índice de dominância), conectividade estrutural (segregação intra/inter-comunidade), e hierarquia (classificação núcleo/intermediária/periférica). Visualizações específicas incluem heatmaps de rivalidades estruturais, scatter plots relacionando tamanho de comunidade com PageRank médio, e distribuições de índices de dominância.

O dashboard implementa controles de filtragem que permitem seleção dinâmica de esporte, gênero e período temporal, atualizando instantaneamente todas as visualizações. Esta interatividade facilita análises exploratórias e identificação de padrões específicos de contextos competitivos particulares. Embora o dashboard não constitua contribuição metodológica central deste trabalho, sua disponibilização como artefato complementar demonstra aplicabilidade prática da abordagem e oferece ferramenta acessível para exploração dos resultados por audiências não-técnicas.